\chapter{Supplementary Optimization Diagnostics} \label{chap:appendix-b}

This appendix records the later log-loss-focused modeling pass that followed the main predictive analysis in Chapter~\ref{chap:protein_results}. The goal of that pass was narrow: train an attention model whose checkpoint selection explicitly targets validation log loss and then test whether a three-seed ensemble of such models should replace the original frozen ensemble. The resulting comparison is useful because it separates two goals that are related but not identical in this thesis: calibrated probability loss on the one hand, and ranking- and planning-oriented surrogate quality on the other.

The frozen surrogate used for inference in Chapter~\ref{chap:protein_results} is a seed-averaged ensemble, so it does not have a single epoch-by-epoch optimization path. For that reason, the optimization curve here is shown for the log-loss-optimized single-attention run, while the ensemble-level comparison is summarized in tables. Figure~\ref{fig:appb_training_loss_curve} tracks train and validation loss across epochs, and Table~\ref{tab:appb_loss_summary} reports the final split-level loss summary at the selected checkpoint. The held-out test split is summarized only once at the selected checkpoint and is not traced across epochs, so that it remains a reporting set rather than a tuning signal.

\begin{figure}[ht]
\centering
\includegraphics[width=0.92\textwidth]{appendix_b_training_loss_curve.png}
\caption[Optimization curve for the log-loss attention run]{Optimization trajectory for the log-loss-optimized sequence-plus-text ESM2-3B attention run used in this appendix. The figure shows train loss and validation loss over epochs, with the selected checkpoint marked by the dashed vertical line. The held-out test split is not evaluated epoch by epoch.}
\label{fig:appb_training_loss_curve}
\end{figure}

The main use of Figure~\ref{fig:appb_training_loss_curve} is diagnostic rather than comparative. It shows whether the log-loss-optimized attention model reaches a stable region before early stopping and whether the validation trajectory supports the selected checkpoint. In this appendix, the validation curve is the relevant companion to the training curve because model selection is based on validation behavior, not on the held-out test split.

\begin{table}[ht]
\centering
\small
\begin{tabularx}{0.82\textwidth}{@{}>{\raggedright\arraybackslash}X*{3}{>{\centering\arraybackslash}X}@{}}
\toprule
Split & Log loss & Brier & PR-AUC \\
\midrule
Intra train & 0.653 & 0.230 & 0.675 \\
Intra val & 0.755 & 0.257 & 0.664 \\
Intra test & 0.748 & 0.249 & 0.682 \\
BIOGRID test & 0.572 & 0.193 & 0.790 \\
\bottomrule
\end{tabularx}
\caption[Final split-level loss summary for the log-loss attention run]{Final split-level loss and ranking summary for the log-loss-optimized single-attention run used in Appendix~\ref{chap:appendix-b}. The held-out test split is summarized only at the selected checkpoint rather than traced across epochs.}
\label{tab:appb_loss_summary}
\end{table}


Table~\ref{tab:appb_loss_summary} complements the optimization curve by giving the final split-level loss summary for the same log-loss-optimized run. Reporting the held-out test loss only once, after checkpoint selection, preserves the logic of the evaluation design used throughout the thesis.

\begin{table}[ht]
\centering
\scriptsize
\begin{tabularx}{\textwidth}{@{}>{\raggedright\arraybackslash}X*{8}{>{\centering\arraybackslash}X}@{}}
\toprule
Split & \multicolumn{4}{c}{Old} & \multicolumn{4}{c}{New} \\
\cmidrule(lr){2-5}\cmidrule(lr){6-9}
& log loss & Brier & PR & rho2 & log loss & Brier & PR & rho2 \\
\midrule
BIOGRID test & 0.795 & 0.238 & 0.788 & 0.238 & 0.572 & 0.193 & 0.790 & 0.273 \\
Intra test & 0.890 & 0.261 & 0.697 & 0.124 & 0.748 & 0.249 & 0.682 & 0.107 \\
Intra val & 0.940 & 0.277 & 0.687 & 0.099 & 0.755 & 0.257 & 0.664 & 0.073 \\
\bottomrule
\end{tabularx}
\caption[Single-run tradeoff under log-loss optimization]{Comparison between the earlier representative single-attention run and the later log-loss-optimized single-attention run. Lower log loss and Brier indicate better calibration-oriented performance, whereas PR-AUC and pilot rho2 summarize ranking and planning-oriented quality.}
\label{tab:appb_single_tradeoff}
\end{table}


Table~\ref{tab:appb_single_tradeoff} compares the earlier representative single-attention run used in the main predictive table with the later log-loss-optimized single-attention run. The direction of the tradeoff is clear. Optimizing directly for log loss lowers both log loss and Brier on all three key splits, and it slightly improves the external BIOGRID PR-AUC. At the same time, it reduces the in-domain PR-AUC and lowers the pilot \(\rho^2\) values that feed the planning interpretation in Chapter~\ref{chap:protein_results}. In other words, the later run is better calibrated in the loss-based sense, but it is not automatically a stronger surrogate for the later ranking- and planning-based uses of the model.

\begin{table}[ht]
\centering
\scriptsize
\begin{tabularx}{\textwidth}{@{}>{\raggedright\arraybackslash}X*{8}{>{\centering\arraybackslash}X}@{}}
\toprule
Split & \makecell{Old\\log loss} & \makecell{New\\log loss} & \makecell{Old\\Brier} & \makecell{New\\Brier} & \makecell{Old\\PR} & \makecell{New\\PR} & \makecell{Old\\rho2} & \makecell{New\\rho2} \\
\midrule
BIOGRID test & 0.685 & 0.552 & 0.218 & 0.187 & 0.799 & 0.802 & 0.265 & 0.295 \\
Intra test & 0.784 & 0.722 & 0.242 & 0.248 & 0.715 & 0.649 & 0.143 & 0.084 \\
Intra val & 0.834 & 0.731 & 0.261 & 0.255 & 0.697 & 0.641 & 0.106 & 0.059 \\
\bottomrule
\end{tabularx}
\caption[Ensemble tradeoff under log-loss optimization]{Comparison between the original frozen attention ensemble and the log-loss-optimized attention ensemble under a common rescoring path. These are the quantities used in the balanced replacement rule discussed in Appendix~\ref{chap:appendix-b}.}
\label{tab:appb_ensemble_tradeoff}
\end{table}


Table~\ref{tab:appb_ensemble_tradeoff} carries the same comparison to the ensemble level. The three-seed log-loss ensemble improves log loss on all three key splits and improves Brier on two of the three. However, it fails the balanced replacement rule adopted for the thesis: the in-domain test Brier score worsens slightly, the in-domain PR-AUC drops from 0.715 to 0.649, and the pilot \(\rho^2\) on the Intra validation split falls from 0.106 to 0.059. The log-loss-optimized ensemble is therefore retained only as a calibration-oriented contrast, while the original frozen attention ensemble remains the main surrogate used in Chapter~\ref{chap:protein_results}. This is the substantive reason the thesis does not select models by log loss alone.
